\section*{第 \ref{ch:b1:nucleophilic-substitution} 章——亲核取代}

\begin{solution}{exo:b1:nucleophilic-substitution:1}
\ce{CH3CH2OH + Br-}; \ce{CH3OCH2CH3 + I-}; \ce{CH3CH2CH2CN + Cl-};
\ce{CH3NH3+ + Br-}。\omterm{def:b1:nucleophilic-substitution:substitution}{底物}：卤代烷；\omterm{def:b1:electronic-effects:nucleophile}{亲核试剂}：
\ce{OH-}、\ce{CH3CH2O-}、\ce{CN-}、\ce{NH3}；\omterm{def:b1:electronic-effects:nucleophile}{离去基团}：卤离子。
\end{solution}

\begin{solution}{exo:b1:nucleophilic-substitution:2}
二级：速率变为三倍；两者均减半时速率除以 4。仅对 RY 为一级：
不变；减半。
\end{solution}

\begin{solution}{exo:b1:nucleophilic-substitution:3}
SN2（伯，良好\omterm{def:b1:electronic-effects:nucleophile}{亲核试剂}，\omterm{def:b1:intermolecular-forces-solvents:solvent-classes}{非质子性溶剂}）；SN1（叔，水）；SN2
（甲基，强\omterm{def:b1:electronic-effects:nucleophile}{亲核试剂}）；SN1 \omterm{def:b1:nucleophilic-substitution:sn1}{溶剂解}（仲，弱中性
\omterm{def:b1:electronic-effects:nucleophile}{亲核试剂}，\omterm{def:b1:intermolecular-forces-solvents:solvent-classes}{质子性溶剂}），可能伴有部分 SN2。
\end{solution}

\begin{solution}{exo:b1:nucleophilic-substitution:4}
氰根从与溴相对的位置进入，其余三个
基团翻转。产物中，CN 基团与原来的 Br 一样排第一，其余
基团的排序不变：翻转后的\omterm{def:b1:stereochemistry-in-depth:isomers}{构型}具有相反的
\omterm{def:b1:stereochemistry-in-depth:isomers}{构型}标记，为 \cip{S}-2-甲基丁腈。
\end{solution}

\begin{solution}{exo:b1:nucleophilic-substitution:5}
1. \omterm{def:b1:electronic-effects:cleavage}{异裂}：箭头从 C--Br 键指向 Br，生成 \ce{(CH3)3C+} 和
\ce{Br-}（慢步）。2. 水的\omterm{def:b1:lewis-resonance-vsepr:lewis}{孤电子对}进攻带正电的碳：
\ce{(CH3)3C-OH2+}。3. 另一个水分子的\omterm{def:b1:lewis-resonance-vsepr:lewis}{孤电子对}夺取
\ce{OH2+} 基团上的一个质子，O--H 键的电子对回到氧原子上。只有第一步
慢步决定\omterm{def:b1:rate-laws:rate}{反应速率}；作为溶剂的水大量过量，
其浓度基本不变。
\end{solution}

\begin{solution}{exo:b1:nucleophilic-substitution:6}
与 Br 相连的碳是\omterm{def:b1:electronic-effects:carbon-class}{伯碳}，但与其相邻的叔丁基
占据了 C--Br 键背后的空间：SN2 \omterm{def:b1:elementary-steps:energy-profile}{过渡态}十分拥挤。
SN1 则需生成伯\omterm{def:b1:electronic-effects:intermediates}{碳正离子}，而它过于不稳定。
\end{solution}

\begin{solution}{exo:b1:nucleophilic-substitution:7}
碘离子通过 SN2 取代碘离子，每次都伴随翻转：\cip{S} 变为
\cip{R}，直至两者等量存在，即\omterm{def:b1:stereochemistry-in-depth:enantiomers}{外消旋混合物}。每一次
取代都使一个 \cip{S} 分子变为一个 \cip{R} 分子，旋光度
减少两个分子的量，因此旋光度衰减
速率为\omterm{def:b1:nucleophilic-substitution:substitution}{底物}被取代速率的两倍。
\end{solution}

\begin{solution}{exo:b1:nucleophilic-substitution:8}
SN2：溴甲烷 $>$ 1-溴丙烷 $>$ 2-溴丙烷 $>$
2-溴-2-甲基丙烷（空间拥挤）。SN1：顺序相反（\omterm{def:b1:electronic-effects:intermediates}{碳正离子}
稳定性），溴甲烷和 1-溴丙烷几乎不反应。
\end{solution}

\begin{solution}{exo:b1:nucleophilic-substitution:9}
$x_{\text{翻转}} - x_{\text{保留}} = 0.12$ 且 $x_{\text{翻转}} + x_{\text{保留}} = 1$：
\qty{56}{\percent} 翻转，\qty{44}{\percent} 保持\omterm{def:b1:stereochemistry-in-depth:isomers}{构型}。以 SN1
\omterm{def:b1:nucleophilic-substitution:sn1}{外消旋化}为主，伴有略多的翻转：离去的溴离子
水到达时仍遮挡其一侧。
\end{solution}

\begin{solution}{exo:b1:nucleophilic-substitution:10}
$v = k_1k_2[\mathrm{RY}][\mathrm{Nu}]/(k_{-1}[\mathrm{Y^-}] + k_2[\mathrm{Nu}])$
（见\cref{prop:b1:nucleophilic-substitution:sn1-law}）。加入 \ce{Y-}
使分母增大：更多\omterm{def:b1:electronic-effects:intermediates}{碳正离子}回到\omterm{def:b1:nucleophilic-substitution:substitution}{底物}而非
被捕获，速率随之下降。当
$k_2[\mathrm{Nu}] \gg k_{-1}[\mathrm{Y^-}]$ 时此项可忽略，
例如\omterm{def:b1:electronic-effects:nucleophile}{亲核试剂}即为溶剂。
\end{solution}

\begin{solution}{exo:b1:nucleophilic-substitution:11}
浓度相等时：$-\mathrm{d}c/\mathrm{d}t = kc^2$，$1/c = 1/C_0 + kt$，
$t_{1/2} = 1/(kC_0) = \qty{5.0e5}{s}$（约六天）。当氢氧根
过量五十倍时，$[\ce{OH-}] \approx 50C_0$ 为常数，$c = C_0
e^{-k't}$，其中 $k' = 50kC_0 = \qty{1.0e-4}{s^{-1}}$，$t_{1/2} = \ln 2/k' =
\qty{6.9e3}{s}$，约两小时。
\end{solution}

\begin{solution}{exo:b1:nucleophilic-substitution:12}
溴离子须驱离 \ce{OH-}，后者是\omterm{def:b1:predominant-reaction:strong-acid}{强碱}，又是
较差的\omterm{def:b1:electronic-effects:nucleophile}{离去基团}，故不反应。在浓 HBr 中，羟基被质子化为
\ce{-OH2+}，其\omterm{def:b1:electronic-effects:nucleophile}{离去基团}为水。叔醇：水离去，
生成\omterm{def:b1:electronic-effects:intermediates}{碳正离子}，再被 \ce{Br-} 捕获（SN1）。伯醇：
\ce{Br-} 从背面进攻 \ce{R-OH2+} 的碳原子，水同时离去
（SN2）。
\end{solution}

\begin{solution}{pb:b1:nucleophilic-substitution:1}
\textbf{1.} 每个参与反应的分子生成一个 \ce{H+} 和一个 \ce{Cl-}；
电导率为各离子贡献之和，与
离子的物质的量成正比。
\textbf{2.} $[\mathrm{RCl}]/[\mathrm{RCl}]_0 = (\kappa_\infty -
\kappa)/\kappa_\infty$.
\textbf{3.} $\ln[\mathrm{RCl}] = \ln[\mathrm{RCl}]_0 - kt$：以
$\ln(\kappa_\infty - \kappa)$ 对 $t$ 作图。
\textbf{4.} 0.693, 0.603, 0.513, 0.423, 0.333, 0.153, $-0.117$, $-0.387$.
\textbf{5.} 是，斜率始终为 $-0.0180$ 每分钟，即一级反应。
\textbf{6.} 对水而言，任意级数均可：其浓度恒定
（级数简并）。
\textbf{7.} $k = 0.0180/60 = \qty{3.0e-4}{s^{-1}}$.
\textbf{8.} $t_{1/2} = \ln 2/k = \qty{2.3e3}{s}$, \qty{38.5}{min}.
\textbf{9.} $\ln 10/k = \qty{7.7e3}{s}$, \qty{128}{min}.
\textbf{10.} $\ln(\kappa_\infty - \kappa)$：$t = 0$ 时为 0.693，0.109 对应
\qty{10}{min}，依此类推：斜率为 $-0.0584$ 每分钟，$k = \qty{9.7e-4}{s^{-1}}$。
\textbf{11.} $9.7/3.0 = 3.2$.
\textbf{12.} 在 \qty{30}{min} 时：$2.000(1 - e^{-0.0584 \times 30}) = 1.653$，
实测值为 1.654。
\textbf{13.} SN1：叔\omterm{def:b1:nucleophilic-substitution:substitution}{底物}，\omterm{def:b1:intermolecular-forces-solvents:solvent-classes}{质子性溶剂}，中性\omterm{def:b1:electronic-effects:nucleophile}{亲核试剂}。
\textbf{14.} 与\cref{exo:b1:nucleophilic-substitution:5}相同，改用 Cl。
\textbf{15.} 慢的电离步骤只涉及\omterm{def:b1:nucleophilic-substitution:substitution}{底物}：$v =
k[\mathrm{RCl}]$.
\textbf{16.} 氢氧根只能在快速的捕获步骤中起作用，即在
决速电离步骤之后；速率与之无关。
\textbf{17.} \omterm{def:b1:electronic-effects:intermediates}{碳正离子}为平面结构且非\omterm{def:b1:stereochemistry-in-depth:stereocentre}{手性}，水从其两个面
等概率进攻，得到外消旋醇。
\textbf{18.} 生成\omterm{def:b1:electronic-effects:intermediates}{碳正离子}的第一个势垒较高（电离步），\omterm{def:b1:electronic-effects:intermediates}{碳正离子}
位于较浅的势阱中，随后捕获的势垒较低：第一步
为\omterm{def:b1:elementary-steps:rds}{决速步}。
\textbf{19.} 将生成伯\omterm{def:b1:electronic-effects:intermediates}{碳正离子}（过于不稳定），而且水
作为\omterm{def:b1:electronic-effects:nucleophile}{亲核试剂}过弱，难以快速进行 SN2 反应。
\textbf{20.} $k = A\,e^{-E_a/RT}$.
\textbf{21.} $E_a = R\ln(k_2/k_1)/(1/T_1 - 1/T_2)$.
\textbf{22.} $E_a = 8.314 \times \ln(9.74/3.00)/(1/298.15 - 1/308.15) =
\qty{9.0e4}{J/mol}$.
\textbf{23.} 电离步骤的势垒高度，即
\omterm{def:b1:elementary-steps:rds}{决速步}。
\textbf{24.} $k = 3.0 \times 10^{-4} \exp\bigl(\frac{90\,000}{8.314}
(\frac{1}{298.15} - \frac{1}{318.15})\bigr) = \qty{2.9e-3}{s^{-1}}$.
\textbf{25.} \textbf{$E_a = \qty{90}{kJ/mol}$}.
\end{solution}
